\section{能量-动量张量}\label{sec:emt}

在连续的 $D$ 维时空中，\emt\ 的规范不变部分为
\begin{equation}
  \begin{split}
    T_{\mu\nu}(x) \equiv \frac{1}{g_0^2}\left[
      \calo_{1,\mu\nu}(x)-\frac{1}{4}\calo_{2,\mu\nu}(x)\right]
    + \frac{1}{4} \calo_{3,\mu\nu}(x)
    - \frac{1}{2} \calo_{4,\mu\nu}(x)
    - \calo_{5 ,\mu\nu}(x) \,,
    \label{eq:emt}
  \end{split}
\end{equation}
其中 $g_0$ 是 QCD 的裸耦合常数。各算符定义为
\begin{equation}
  \begin{split}
    \calo_{1,\mu\nu}(x) &\equiv
    F_{\mu\rho}^a(x)F_{\nu\rho}^a(x)\,,\\[10pt] \calo_{2,\mu\nu}(x)
    &\equiv
    \delta_{\mu\nu}F_{\rho\sigma}^a(x)F_{\rho\sigma}^a(x)\,,\\[10pt]
    \calo_{3_f,\mu\nu}(x) &\equiv
    \bar\psi_{f}(x)\left(\gamma_\mu\overleftrightarrow{D}\!_\nu
    +\gamma_\nu\overleftrightarrow{D}\!_\mu\right)\psi_{f}(x)\,,\\
    \calo_{4_f,\mu\nu}(x)
    &\equiv \delta_{\mu\nu}\bar\psi_{f}(x)
    \overleftrightarrow{\slashed{D}}\psi_{f}(x)\,,\phantom{\bigg(\bigg)}\\
    \calo_{5_f,\mu\nu}(x)
    &\equiv \delta_{\mu\nu}
    m_{f,0}\bar\psi_{f}(x)\psi_{f}(x)\,,\\
    \calo_{i,\mu\nu}(x) &=
    \sum_{f=1}^{n_F}\calo_{i_f,\mu\nu}(x)\,,\qquad i\in\{3,4,5\}\,,
    \label{eq:ops}
  \end{split}
\end{equation}
其中 $f$ 标记 $n_F$ 种不同的夸克味，$m_{f,0}$ 是裸夸克质量，且
\begin{equation}
  \begin{split}
    F^a_{\mu\nu} = \partial_\mu A_\nu^a - \partial_\nu A_\mu^a +
    f^{abc}A_\mu^bA_\nu^c\,,\qquad
    \overleftrightarrow{D}\!_\mu = \partial_\mu -
    \overleftarrow{\partial}\!_\mu + 2A_\mu\,.
  \end{split}
\end{equation}
用 $i_f\in\{i_1,\ldots,i_{n_F}\}$ 这一记号标记不同味道的下标，在后文中将会
有用。一般地，$T_{\mu\nu}$ 可以含有规范依赖的算符，它们在计算物理矩阵元时
消失~\cite{Nielsen:1977sy}。此处及下文，我们默认所有复合算符的真空期望值均
已被减除（换言之，例如 $\calo_{1,\mu\nu}$ 的精确定义应为
$F_{\mu\rho}^aF_{\rho\nu}^a-\langle F_{\mu\rho}^aF_{\rho\nu}^a\rangle$），
使得 $\langle \calo_{i,\mu\nu}(x)\rangle\equiv 0\ \forall i$。

本文聚焦于考虑 \emt\ 自身的物理矩阵元的情形，即在同时空点上没有其他算符与
\emt\ 相乘。此时，运动方程（\eom{}）使 \eqn{eq:ops} 中的算符组变得冗余。
特别地，常规 QCD 中夸克场的 \eom\ 给出
\begin{equation}
  0 = \calo_{4, \mu \nu} (x) + 2 \, \calo_{5, \mu \nu} (x) \,,
  \label{eq:eom}
\end{equation}
这使我们能从 \eqn{eq:ops} 的算符组中消去 $\calo_{5,\mu\nu}$。注意由于该
关系，\eqn{eq:emt} 中最后两项相消。

我们进一步假定所有夸克质量彼此相等：
\begin{equation}
	m_{f, 0} = m_0\,, \qquad f=1,\ldots,n_F\,.
\end{equation}
于是不同夸克不可分辨，两种不同夸克味之间的混合不能依赖于味道。

对流动场定义与 \eqn{eq:ops} 类似的算符，我们写作
\begin{equation}
  \begin{split}
    \tcalo_{1,\mu\nu}(t,x) &\equiv
    G_{\mu\rho}^a(t,x)G_{\nu\rho}^a(t,x)\,,\\[10pt]
    \tcalo_{2,\mu\nu}(t,x) &\equiv
    \delta_{\mu\nu}G_{\rho\sigma}^a(t,x)G_{\rho\sigma}^a(t,x)\,,\\[10pt]
    \tcalo_{3_f,\mu\nu}(t,x) &\equiv Z_{\chi_f}
    \bar\chi_{f}(t,x)\left(\gamma_\mu\overleftrightarrow{\mathcal{D}}\!_\nu
    +\gamma_\nu\overleftrightarrow{\mathcal{D}}\!_\mu\right)\chi_{f}(t,x)\,,\\
    \tcalo_{4_f,\mu\nu}(t,x)
    &\equiv Z_{\chi_f} \delta_{\mu\nu}\bar\chi_{f}(t,x)
    \overleftrightarrow{\slashed{\mathcal{D}}}\chi_{f}(t,x)\,,
    \\
    \tcalo_{i,\mu\nu} &= \sum_{f=1}^{n_F}\tcalo_{i_f,\mu\nu}\,,\qquad i\in\{3,4\}\,,
    \label{eq:tops}
  \end{split}
\end{equation}
其中 $Z_{\chi_f}\equiv Z_\chi$ 是流动夸克场的重正常数，而
\begin{equation}
  \begin{split}
    \overleftrightarrow{\mathcal{D}}\!_\mu = \partial_\mu -
    \overleftarrow{\partial}\!_\mu + 2B_\mu\,.
    \label{eq:leftrightcalD}
  \end{split}
\end{equation}
既然已用 \eqn{eq:eom} 把 $\calo_{5,\mu\nu}$ 从算符组中消去，就不必在
\eqn{eq:tops} 中包含该算符的流动版本。与常规 QCD 的复合算符类似，我们假定
流动复合算符的真空期望值已被减除，
即 $\langle\tcalo_{i,\mu\nu}(t,x)\rangle\equiv 0\ \forall i$。

现在可以利用小流时间下的展开\,\cite{Luscher:2011bx}
\begin{equation}
  \begin{split}
    \tcalo_{i,\mu\nu}(t,x) = \zeta_{ij}(t)\calo_{j,\mu\nu}(x) + \ldots\,,
    \label{eq:zetas}
  \end{split}
\end{equation}
来建立流动算符与常规 QCD 算符之间的关系。在 \eqn{eq:zetas} 以及下文中，
约定对 $j$ 从 $1$ 到 $4$ 求和。省略号代表当 $t\to 0$ 时趋于零的项，本文通篇
将其忽略。如上所述，该方程左边的矩阵元在 QCD 参数重正化之后是有限的，而右边
常规 QCD 算符的矩阵元一般是发散的。因此混合矩阵 $\zeta_{ij}(t)$ 也是发散的。

将 \eqn{eq:zetas} 求逆，并用它把能动量张量中的常规 QCD 算符重新用流动场
表达，便得到
\begin{equation}
  \begin{split}
    T_{\mu\nu}(x) =  c_i (t) \tcalo_{i, \mu \nu} (t, x) \,,
    \label{eq:emtflow}
  \end{split}
\end{equation}
其中
\begin{equation}
  \begin{split}
    c_i(t) \equiv \frac{1}{g_{0}^2} \left( \zeta_{1i}^{-1} (t)
		 - \frac{1}{4} \zeta_{2i}^{-1} (t) \right)
		 + \frac{1}{4} \zeta_{3i}^{-1} (t)
                         \,,\qquad i=1,\ldots,4\,.
    \label{eq:EMTGF}
	\end{split}
\end{equation}
由于 $\tcalo_i$ 的矩阵元以及能动量张量本身都是有限的（在质量与电荷重正化之后），\eqn{eq:EMTGF} 中的普适系数 $c_i(t)$ 也是有限的。\citere{Makino:2014taa}已在微扰论中把它们算到 \nlo\ QCD。本文的目标是把它们求值到 \nnlo\ QCD。
